\begin{remark}\label{rem:Q3}
At $Q^2=3$, where $\hat f(3)=1122$, positivity holds for $D\le\Dexact$ by Theorem~\ref{thm:bounded} and for all $D$ under Hypothesis~\ref{hyp:R} by Theorem~\ref{thm:all}. An unconditional proof along the lines of Sections~\ref{sec:Q1} and~\ref{sec:Q2} would need the closed form of $\varphi_3$. The coefficients of the product formula that we computed (all $x^ey^\ell$ with $e<60$) are matched by $g(\frac\sigma2)\varphi_3=1122\,\varphi^{\rm opt}_{2,3}+J_3=1122\,\mathcal A_{2,3}-1122\,\cH|V_3+J_3$, with $\varphi^{\rm opt}_{2,3}$ the form of \cite[(9.12)--(9.13)]{DMZ12} and $J_3$ a holomorphic Jacobi form of weight two and index three built from $A$, $B$ and Eisenstein series at $\sigma$ and $\sigma/2$. Proving this identity needs the analogue of Proposition~\ref{prop:P3} at index four, and the signs of the theta components of $J_3$ would then have to be determined.
\end{remark}

\section{What remains}\label{sec:remains}

The proof of Theorem~\ref{thm:all} uses Hypothesis~\ref{hyp:R} and nothing else that is unproved. Three routes would remove it.
\begin{enumerate}
\item[(1)] \emph{The expansion (H).} The method of Bringmann and Manschot \cite{BMa13} writes a mixed mock Jacobi form as a combination of Poincar\'e series and of the image of the shadow under a non-holomorphic Eichler integral, and proves that the difference is a cusp form of positive weight that vanishes for weight reasons. At level two the Poincar\'e series have to be summed over the three frames, and the vanishing argument has to be made for $\Gamma^0(2)$; the space of holomorphic cusp forms that could intervene has weight $6+\frac12$ in each theta component, and its vanishing is the one point that needs an idea. The agreement to $21$ digits at $Q^2=1,2$ (Section~\ref{ssec:evidence}) shows that no such cusp form is present there.
\item[(2)] \emph{The polar data (P0), (P$\infty$).} Cardoso, Nampuri and Rossell\'o derive the continued-fraction formula \cite[(2.43)]{CNR21} by tracking the poles of $1/\Phit_k$ across the walls of marginal stability. The same contour argument, applied to $\varphi^F_m$ transformed by $\sigma\mapsto-1/\sigma$, would give (P0); (P$\infty$) is its restriction to the walls fixed by the half shift.
\item[(3)] \emph{Exact computation.} Each further unit of $Q^2$ for which the remainder $J_m$ has a closed form, as at $Q^2=1,2$, removes the hypothesis on that row. Theorem~\ref{thm:cert} shows that the margin for $3\le Q^2\le7$ is above $0.53$ as soon as $D>500$, so on these rows the analytic picture of Section~\ref{sec:model} is already in force at moderate discriminant.
\end{enumerate}
The method itself is not special to $N=2$. For the $\ZZ_N$ CHL models with $N=3,5,7$ the leading family of walls again resums, now to $|1/\eta(\tau)^{k+2}\eta(N\tau)^{k+2}|^2$ evaluated at $e^{-2\pi x/N+i\theta}$ with $k+2=24/(N+1)$, and positivity reduces in the same way to a lower bound for this modulus and to an expansion of type (H) at level $N$.
